\documentclass[11pt]{article}
\usepackage[a4paper,margin=26mm,headheight=15pt]{geometry}
\usepackage[T1]{fontenc}
\usepackage{lmodern,microtype,amsmath,amssymb,amsthm,mathtools,mathrsfs}
\usepackage{booktabs,longtable,array,enumitem,xcolor,xurl,fancyhdr}
\usepackage[colorlinks=true,linkcolor=blue!40!black,citecolor=blue!40!black,urlcolor=blue!40!black]{hyperref}
\usepackage{bookmark,needspace}
\numberwithin{equation}{section}
\newtheorem{theorem}{Theorem}[section]
\newtheorem{lemma}[theorem]{Lemma}
\newtheorem{proposition}[theorem]{Proposition}
\newtheorem{corollary}[theorem]{Corollary}
\theoremstyle{definition}\newtheorem{definition}[theorem]{Definition}
\theoremstyle{remark}\newtheorem{remark}[theorem]{Remark}\newtheorem{example}[theorem]{Example}
\newcommand{\C}{\mathbb C}\newcommand{\N}{\mathbb N}\newcommand{\Z}{\mathbb Z}
\newcommand{\D}{\mathscr D}\newcommand{\calH}{\mathcal H}\newcommand{\calR}{\mathcal R}
\newcommand{\calJ}{\mathcal J}\newcommand{\calZ}{\mathcal Z}
\newcommand{\dd}{\,\mathrm d}\newcommand{\Rea}{\operatorname{Re}}
\DeclareMathOperator{\Li}{Li}\DeclareMathOperator{\FP}{FP}\DeclareMathOperator{\Res}{Res}
\newcommand{\coeff}[2]{[#1]#2}
\newcommand{\pin}{\texttt{3fdd6cc447fac0731801588723109ec3b2e21cec}}
\setlength{\parskip}{4pt}\setlength{\parindent}{15pt}\setlength{\emergencystretch}{3em}
\allowdisplaybreaks[2]
\pagestyle{fancy}\fancyhf{}
\fancyhead[L]{\small Arithmetic spectral jets}\fancyhead[R]{\small ProveIt continuation}
\fancyfoot[C]{\thepage}\renewcommand{\headrulewidth}{.3pt}
\hypersetup{pdftitle={Arithmetic Spectral Jets: Stieltjes Residues, Gamma Reciprocity, and Polylogarithmic Untwisting},pdfauthor={Research continuation prepared for Vladimir Reshetnikov},pdfsubject={Exact identities for arithmetic multiplicities and higher-pole spectral finite parts}}
\begin{document}
\begin{titlepage}
\centering
{\small PROVEIT / ANALYSIS / POLYLOGARITHMS\par}
\vspace{18mm}
{\LARGE\bfseries Arithmetic Spectral Jets\par}
\vspace{6mm}
{\Large Stieltjes residues, Gamma reciprocity,\par
and polylogarithmic untwisting\par}
\vspace{13mm}
{\large Research continuation prepared for Vladimir Reshetnikov\par}
\vspace{3mm}
{\small Mathematical derivations and computational checks assisted by ChatGPT\par}
\vspace{5mm}
{10 October 2026\par}
\vspace{10mm}
\begin{abstract}
\noindent
The all-order local calculus for polynomial spectral multiplicities extends to arithmetic
weights, but ordinary derivatives at the origin generally cease to exist. We prove a
finite-pole transfer theorem that expresses every normalized Laurent jet as a polynomial
in the factor exponents. Its terms beyond the ordinary Taylor degree are determined
entirely by the principal parts of the base Dirichlet series. This answers the explicitly
posed arithmetic-multiplicity extension for quasi-polynomial weights and for products
of Hurwitz zeta functions, including divisor multiplicities and several distinct poles.

For multiplicative Hurwitz lattices, all local coefficients are finite polynomials in
generalized Stieltjes constants. Alternation gives exact algebraic identities between
Laurent jets: a six-term first-jet identity for double-divisor multiplicities equals a
Vandermonde divided by nine. A convergent centered Gamma series is symmetric in two
independent Hurwitz parameters; its derivatives and anchored primitive yield reciprocal
polygamma and negative-zeta-derivative identities. Generalized harmonic numbers describe
every higher single-shift jet. Polylogarithmic twists lower the pole order, and a finite,
explicit Gamma--Stieltjes compensation restores commutation of Abel untwisting with
Laurent coefficient extraction, at every order and for every polynomial multiplicity.

All theorem statements have ordinary analytic or algebraic proofs. Exact finite
regression tests and independent high-precision diagnostics are supplied separately.
The first multiplicative anomaly and the classical special-function ingredients retain
their attribution. No claim of worldwide priority, arithmetic independence, or
proof-assistant formalization is made.
\end{abstract}
\vfill
\begin{minipage}{.94\textwidth}\small\raggedright
\textbf{Source anchor.} Targeted repository inspection is pinned to\
\pin.\par
\textbf{Scope.} The report addresses question Q10 in the existing
\emph{Resonant Gamma Jets and Directional Zeta Identities} continuation. It does not settle the Gaussian
$S_6$ conjecture or the distinct current $S_8$ conjecture.
\end{minipage}
\end{titlepage}
\begingroup
\small
\setlength{\parskip}{0pt}
\tableofcontents
\endgroup
\clearpage

% BEGIN sections/01-scope.tex
\section{Research target, contribution, and conventions}
\label{asj:sec:scope}

\subsection{A specific repository question}
The manuscript's Hurwitz-jet chapters organize parameter derivatives, spectral
derivatives, generalized Stieltjes constants, and negative-order polygamma in one
calculus \cite{repo}. A later report already proves a normalized-exponent polynomial
law for polynomial spectral multiplicities, including a polynomial times a power of
a logarithm \cite{prior}. Its question Q10 asks for arithmetic and nonpolynomial
multiplicities, beginning with periodic or quasi-polynomial weights and then moving
to base zeta functions with several or higher-order poles.

That is the precise target here. The extension is not obtained by assuming that an
arbitrary arithmetic Dirichlet series continues meromorphically. Instead, we state
such a continuation as a hypothesis in the general theorem, prove the requisite
normal convergence, and verify the hypotheses unconditionally in explicit families.
For these families the result is a complete coefficient algorithm rather than an
asymptotic extrapolation.

The new scope relative to \cite{prior} consists of a general finite-principal-part
transfer theorem, the separation of the mean and oscillatory parts of
quasi-polynomial multiplicities, independent Hurwitz factors with all their
Stieltjes residue layers, explicit Gamma reciprocity and primitive formulas, and
an all-order finite compensation for removing a polylogarithmic twist. The
alternating identities extend the report's simple-pole alternants to actual
arithmetic higher poles. The argument is elementary once the correct germ is
retained; its simplicity is not a reason to identify a Laurent coefficient with
an undefined derivative.

\subsection{Classical inputs and the limits of a novelty assertion}
The Hurwitz continuation, its parameter derivative, and its relation to
$\log\Gamma$ are standard \cite{dlmf-hurwitz}. The Pochhammer expansion, the
Gamma product, and the polylogarithm expansion near unity are also classical
\cite{dlmf-gamma,dlmf-polylog}. First-derivative multiplicative-anomaly formulas
have an extensive prior theory; Dowker's treatment includes products of linear
factors and polynomial degeneracies \cite{dowker}. We use that theory as context,
not as a new result. The present proofs establish the stated extensions relative
to the inspected repository. The literature inspection does not prove that every
individual consequence is absent from all earlier publications.

The source audit is targeted. The canonical chapters, the relevant earlier
spectral-product section and its research questions, and the incoming-directory
inventory were inspected. The incoming ZIP files were not all extracted and
audited. The artifact records the source revision and the precise question being
answered; it does not assert a line-by-line audit of the entire repository.

\subsection{Laurent jets and local coefficients}
Throughout, $[s^r]F(s)$ denotes a Laurent coefficient. When $r\ge0$, define
\begin{equation}\label{asj:eq:jetdef}
 \calJ_r F=r![s^r]F(s).
\end{equation}
This equals $F^{(r)}(0)$ only when $F$ is holomorphic there. A ``first jet''
means $[s]F(s)$ under this convention, not the finite part of $F'(s)$ under a
separately chosen change of coordinate. The regulator coordinate $s$ is fixed.

For an expansion at infinity, $[x^{-p}]f(x)$ means the coefficient in the
Laurent series in $x^{-1}$. In particular, $[x^{-1}]$ is the negative of the
complex-analytic residue at infinity when that terminology applies. This
coefficient is computed from a finite truncation in all formulas below.

We use the standard Stieltjes normalization
\begin{equation}\label{asj:eq:stieltjes}
 \zeta(1+s,b)=\frac1s+\sum_{n\ge0}\frac{(-1)^n}{n!}\gamma_n(b)s^n,
 \qquad b>0.
\end{equation}
Thus $\gamma_0(b)=-\psi(b)$ and $\gamma_0(1)=\gamma$, Euler's constant.
Spectral superscripts on $\zeta$ differentiate its first argument; the
superscript on $\psi^{(k)}$ is an argument derivative. All logarithms of positive
real quantities are real. The proofs also give local holomorphic shift
continuations on any domain where compatible logarithms have been fixed.

\subsection{What is not being claimed}
An exact polynomial identity in the exponent variables is not an independence
theorem for its special-function coefficients. A finite numerical residual is
not a proof of equality. Neither a new Gamma-series coordinate nor a zeta-jet
coordinate is claimed to be irreducible in a smaller arithmetic algebra.
The ordinary proofs here have not been formalized in Lean or another proof
assistant.

% END sections/01-scope.tex


% BEGIN sections/02-transfer.tex
\section{A transfer theorem for general discrete spectra}
\label{asj:sec:transfer}

\subsection{Hypotheses and a convergent expansion}
Let $(\lambda_\nu)$ be a positive discrete spectrum, counted with multiplicity,
with positive minimum $\lambda_*$ and only finitely many indices below each
fixed bound. Let $c_\nu\in\C$. Assume
\begin{equation}\label{asj:eq:base}
 \D(s)=\sum_\nu c_\nu\lambda_\nu^{-s}
\end{equation}
converges absolutely for $\Rea s>\sigma_a$, for some finite $\sigma_a$, and
extends meromorphically to $\C$. The weights need not be positive. This is a
hypothesis on $\D$, not a claimed property of every arithmetic sequence.

Let
\[
 P(x)=\sum_{k=0}^{d}p_kx^k,
 \qquad a_j> -\lambda_*,\quad 1\le j\le m,
 \qquad U=\sum_{j=1}^m u_j.
\]
Initially in $\Rea U>\sigma_a+d$, define
\begin{equation}\label{asj:eq:Zdef}
 \calZ_P(\boldsymbol u;\boldsymbol a)
 =\sum_\nu c_\nu P(\lambda_\nu)
                   \prod_{j=1}^m(\lambda_\nu+a_j)^{-u_j}.
\end{equation}
The finite shifts do not change the absolute convergence condition at infinity.
Introduce the germs
\begin{align}
 L_j(z)&=\log(1+a_jz),&
 L_{\boldsymbol u}(z)&=\sum_j u_j L_j(z),\label{asj:eq:Ldef}\\
 C_\ell(\boldsymbol u)&=[z^\ell]e^{-L_{\boldsymbol u}(z)}.
 \label{asj:eq:Cdef}
\end{align}
Their coefficients are explicit polynomials:
\begin{equation}\label{asj:eq:Cpoch}
 C_\ell(\boldsymbol u)=
 \sum_{q_1+\cdots+q_m=\ell}
    \prod_{j=1}^m\frac{(-a_j)^{q_j}(u_j)_{q_j}}{q_j!}.
\end{equation}
In particular, $C_0=1$, $C_\ell(0)=0$ for $\ell>0$, and
$\deg C_\ell\le\ell$.

Choose a finite head $E$ consisting of all indices with $\lambda_\nu\le R$,
where $R>\max_j|a_j|$, and put
\[
 \D_E(s)=\D(s)-\sum_{\nu\in E}c_\nu\lambda_\nu^{-s}.
\]
The finite subtraction is entire and leaves all principal parts unchanged.
If the spectrum is finite, the desired conclusions hold with no poles and
no infinite tail; hence suppose the tail is nonempty. Its minimum is denoted
by $M$, so $M>\max|a_j|$.

\begin{lemma}[Normally convergent tail expansion]\label{asj:thm:tail}
The meromorphic continuation of \eqref{asj:eq:Zdef} is
\begin{equation}\label{asj:eq:tail}
 \calZ_P(\boldsymbol u)=
 \calZ_{P,E}(\boldsymbol u)
 +\sum_{k=0}^d p_k\sum_{\ell\ge0}
          C_\ell(\boldsymbol u)\D_E(U+\ell-k),
\end{equation}
where $\calZ_{P,E}$ is the finite original head. On compact subsets away
from the displayed polar divisors, the tail converges normally, as do
all fixed mixed derivatives in $\boldsymbol u$.
\end{lemma}
\begin{proof}
Select $\rho$ with $M^{-1}<\rho<(\max|a_j|)^{-1}$, interpreting the upper
endpoint as infinity when every shift is zero. Cauchy's coefficient estimate
on $|z|=\rho$ gives
\[
 |C_\ell(\boldsymbol u)|\le C_K\rho^{-\ell}
\]
uniformly on every compact exponent set $K$. In a sufficiently far right
half-plane, the absolutely convergent series for $\D_E$ gives, for bounded $v$,
\[
 |\D_E(\ell+v)|\le C M^{-\ell}.
\]
Indeed fix $A>\sigma_a$ and majorize each term by
$|c_\nu|\lambda_\nu^{-A}M^{-\ell+A-\Rea v}$ for large $\ell$.
The resulting geometric factor is $(M\rho)^{-\ell}<1$.
Fixed derivatives of $\D_E$ add powers of $\log\lambda_\nu$; choosing a
slightly larger absolute-convergence exponent absorbs them. Derivatives of
$C_\ell$ are controlled by Cauchy estimates on a larger compact exponent set.

Expanding every shifted factor by its binomial series proves
\eqref{asj:eq:tail} in the initial convergence region, where these estimates
justify interchanging the sums. The normally convergent right side then
supplies the meromorphic continuation. Only finitely many small $\ell$
can contribute a pole on a fixed compact exponent set, because the other
arguments lie in an absolute-convergence half-plane.
\end{proof}

\subsection{The complete local polar numerator}
Let $\mathcal P$ be the set of integer poles $p$ of $\D$ with $p\ge-d$.
It is finite: there are no poles to the right of an absolute-convergence
half-plane. Write the principal part at each such pole as
\begin{equation}\label{asj:eq:principal}
 \D(p+s)=\sum_{h=1}^{H_p}\frac{\kappa_{p,h}}{s^h}
                      +\text{a holomorphic function of }s.
\end{equation}
The highest coefficient $\kappa_{p,H_p}$ is nonzero. Define
\begin{equation}\label{asj:eq:Rdef}
 \calR_p(\boldsymbol u)=
       \sum_{\substack{0\le k\le d\\k+p\ge0}}p_k C_{k+p}(\boldsymbol u).
\end{equation}
A pole with $\calR_p=0$ contributes nothing. Noninteger poles do not pass
through the origin of exponent space.

\begin{theorem}[Finite-principal-part transfer]\label{asj:thm:transfer}
In a neighborhood of $\boldsymbol u=0$, the exact germ is
\begin{equation}\label{asj:eq:germ}
 \boxed{\displaystyle
 \calZ_P(\boldsymbol u)
  =\calH_P(\boldsymbol u)
       +\sum_{p\in\mathcal P}\sum_{h=1}^{H_p}
            \frac{\kappa_{p,h}\calR_p(\boldsymbol u)}{U^h},}
\end{equation}
where $\calH_P$ is holomorphic. The numerator $\calR_p$ is a polynomial
of degree at most $d+p$. If $p>0$ then $\calR_p(0)=0$.
Neither the principal coefficients nor the numerators depend on the
finite head used in Lemma~\ref{asj:thm:tail}.
\end{theorem}
\begin{proof}
A term $\D_E(U+\ell-k)$ has a pole through $U=0$ precisely when
$p=\ell-k$ is an integer pole of $\D$. For a fixed such $p$, its
coefficient is $\sum_k p_k C_{k+p}=\calR_p$. Subtract the finitely many
principal parts \eqref{asj:eq:principal} from these terms. The remaining
small-$\ell$ terms are holomorphic near the origin, and the large-$\ell$
tail is holomorphic there by Lemma~\ref{asj:thm:tail}. Choose the neighborhood
small enough to avoid every other pole of the finitely many remaining
small-$\ell$ factors. This proves \eqref{asj:eq:germ}.
The degree and constant-term assertions follow from \eqref{asj:eq:Cpoch}.
Finally, altering a finite head changes $\D_E$ by an entire function,
which cannot change any principal coefficient. Formula \eqref{asj:eq:Rdef}
uses only $P$ and the shifts.
\end{proof}

A zero numerator at $\boldsymbol u=0$ does not imply a holomorphic ray
restriction when $H_p>1$. It removes at most one order of the pole without
additional cancellation. Conversely, if $\D$ has no relevant integer poles,
then $\calZ_P$ is jointly holomorphic at the origin even if $\D$ has poles
elsewhere. These are local statements at the specified spectral point.

% END sections/02-transfer.tex


% BEGIN sections/03-local.tex
\section{All-order local identities and finite differences}
\label{asj:sec:local}

Normalize the exponent vector by $\sum_j\alpha_j=1$, and set
\begin{equation}\label{asj:eq:Jnormalized}
 \calJ_r(\boldsymbol\alpha)
  =r![s^r]\calZ_P(s\boldsymbol\alpha),\qquad r\ge0.
\end{equation}
Let $F_{[q]}$ denote the homogeneous Taylor part of degree $q$ of a
holomorphic germ or polynomial $F$. Put
$H=\max_{p\in\mathcal P}H_p$, with $H=0$ if $\mathcal P$ is empty.

\begin{theorem}[Polynomiality of every normalized Laurent jet]
\label{asj:thm:polynomial}
On the affine hyperplane $\sum\alpha_j=1$,
\begin{align}
 \calJ_r(\boldsymbol\alpha)
  &=r!\calH_{P,[r]}(\boldsymbol\alpha)
      +r!\sum_{p,h}\kappa_{p,h}
                       \calR_{p,[r+h]}(\boldsymbol\alpha),\label{asj:eq:Jlocal}\\
 \calR_{p,[q]}(\boldsymbol\alpha)
  &=\frac{(-1)^q}{q!}[x^{-p}]P(x)
       \left(\sum_j\alpha_j\log(1+a_j/x)\right)^q.
 \label{asj:eq:Rhom}
\end{align}
Consequently $\deg\calJ_r\le r+H$. A local term of degree $r+h$ can
occur only when $r+h\le d+p$. Every term beyond degree $r$ is determined
by the finite principal parts of $\D$.
\end{theorem}
\begin{proof}
On the normalized ray $U=s$. Extracting the coefficient of $s^r$ in
\eqref{asj:eq:germ} gives \eqref{asj:eq:Jlocal}. The homogeneous part of degree
$q$ of $e^{-L_{\boldsymbol u}(z)}$ is
$(-1)^qL_{\boldsymbol u}(z)^q/q!$. Using \eqref{asj:eq:Rdef} and then
$z=x^{-1}$ proves \eqref{asj:eq:Rhom}. Every factor of the logarithmic
sum starts with $x^{-1}$, proving the sharper degree cutoff.
\end{proof}

For a tangent vector $\boldsymbol v$ with $\sum_jv_j=0$, write
\[
 \Delta_{\boldsymbol v}f(\boldsymbol\alpha)
     =f(\boldsymbol\alpha+\boldsymbol v)-f(\boldsymbol\alpha),
 \qquad V_{\boldsymbol v}(x)=\sum_jv_j\log(x+a_j).
\]
The $\log x$ terms cancel, so $V_{\boldsymbol v}(x)=O(x^{-1})$.

\begin{theorem}[Top local polarization]\label{asj:thm:top}
Suppose $H\ge1$ and let $q=r+H$. For any $q$ tangent vectors,
\begin{align}\label{asj:eq:top}
 &\Delta_{\boldsymbol v_1}\cdots\Delta_{\boldsymbol v_q}
          \calJ_r(\boldsymbol\alpha)\notag\\
 &\quad=(-1)^q r!\sum_{p:H_p=H}\kappa_{p,H}
           [x^{-p}]P(x)\prod_{i=1}^q V_{\boldsymbol v_i}(x).
\end{align}
The same expression equals the corresponding mixed directional
derivative. Every mixed difference of order $q+1$ vanishes.
\end{theorem}
\begin{proof}
The $q$ differentiations annihilate $\calH_{P,[r]}$ and every local
part of degree below $q$. Polarizing the remaining $q$th power in
\eqref{asj:eq:Rhom} supplies exactly $q!\prod_i V_{\boldsymbol v_i}$.
For a polynomial of degree at most $q$, a mixed $q$th finite difference
is its constant mixed $q$th derivative; integrate that derivative over
the unit $q$-cube to see the equality directly. Differences of the
next order vanish by the degree bound.
\end{proof}

The highest pole alone controls \eqref{asj:eq:top}; it does not control
all lower local differences. The following formula retains every layer.

\begin{theorem}[All finite local differences]\label{asj:thm:all-differences}
Let $q>r$, and put $L_{\boldsymbol\alpha}(x)=
\sum_j\alpha_j\log(1+a_j/x)$. Then
\begin{align}\label{asj:eq:all-differences}
 &\Delta_{\boldsymbol v_1}\cdots\Delta_{\boldsymbol v_q}
       \calJ_r(\boldsymbol\alpha)\notag\\
 &=r!\sum_{p,h}\kappa_{p,h}[x^{-p}]P(x)
      [t^{r+h}]e^{-tL_{\boldsymbol\alpha}(x)}
            \prod_{i=1}^q\bigl(e^{-tV_{\boldsymbol v_i}(x)}-1\bigr).
\end{align}
The right side is a finite algebraic computation. Terms with $q>r+h$
vanish automatically.
\end{theorem}
\begin{proof}
The differences kill the degree-$r$ holomorphic part. Apply a difference
to $e^{-tL_{\boldsymbol\alpha}}$: it multiplies the exponential by
$e^{-tV_{\boldsymbol v}}-1$. Repeating and extracting $t^{r+h}$ in
\eqref{asj:eq:Jlocal} gives the assertion. Each of the $q$ factors is
divisible by $t$, proving the last statement. Coefficients in $x^{-1}$
need only be retained through the largest integer $d+p$ occurring.
\end{proof}

\subsection{Normalization, scaling, and the meaning of locality}
For arbitrary weights $\boldsymbol w$ with $W=\sum_jw_j\ne0$,
\begin{equation}\label{asj:eq:scaling}
 r![s^r]\calZ_P(s\boldsymbol w)
   =W^r\calJ_r(\boldsymbol w/W).
\end{equation}
This remains true for Laurent series: substitute $t=Ws$ before extracting
the coefficient. The excluded hyperplane $W=0$ cannot be restored by
blindly taking a limit in these formulas; it lies in the polar divisor
and needs a separately specified restriction.

``Local'' here has a precise meaning: the expression depends on finitely
many principal coefficients of the base Dirichlet series and finitely many
terms of the large-$x$ logarithmic expansion. It is not a claim that the
individual jets have elementary evaluations. Finite-head changes affect
only $\calH_P$, so any identity obtained by more than $r$ tangent differences
is insensitive to finitely many altered spectral terms.

% END sections/03-local.tex


% BEGIN sections/04-arithmetic.tex
\section{Periodic weights and several distinct arithmetic poles}
\label{asj:sec:arithmetic}

\subsection{A complete quasi-polynomial specialization}
Let $q\ge1$ and let $a_0,\ldots,a_K$ be periodic functions on the positive
integers with period $q$. Define
\[
 c_n=\sum_{d=0}^K n^d a_d(n),\qquad
 \mu_d=\frac1q\sum_{j=1}^q a_d(j).
\]
Residue-class splitting in the absolute-convergence region gives
\begin{equation}\label{asj:eq:quasipoly}
 \D(s)=\sum_{d=0}^K q^{d-s}
                  \sum_{j=1}^q a_d(j)\zeta(s-d,j/q).
\end{equation}
This finite formula proves the continuation hypotheses of
Theorem~\ref{asj:thm:transfer}. Its only possible poles are simple, at $p=d+1$,
with residues $\mu_d$.

\begin{corollary}[Only the mean polynomial contributes local defects]
\label{asj:thm:mean}
For quasi-polynomial multiplicities and any polynomial $P$, all normalized
order-$r$ jets have degree at most $r+1$. For $r+1$ tangent vectors,
\begin{align}\label{asj:eq:mean}
 &\Delta_{\boldsymbol v_1}\cdots\Delta_{\boldsymbol v_{r+1}}
       \calJ_r(\boldsymbol\alpha)\notag\\
 &\quad=(-1)^{r+1}r![x^{-1}]P(x)
           \left(\sum_{d=0}^K\mu_dx^d\right)
                       \prod_i V_{\boldsymbol v_i}(x).
\end{align}
If every $\mu_d=0$, the degree is at most $r$ and the displayed difference
vanishes. More generally, replacing $c_n$ by its mean polynomial changes
only a degree-at-most-$r$ part of the normalized jet.
\end{corollary}
\begin{proof}
Use the residues in \eqref{asj:eq:quasipoly} in \eqref{asj:eq:top} and observe
$[x^{-(d+1)}]f(x)=[x^{-1}]x^d f(x)$. Subtracting the mean polynomial
leaves a base Dirichlet series whose possible poles have all been canceled,
so its shifted germ is holomorphic at the exponent origin. Theorem
\ref{asj:thm:polynomial} then gives degree at most $r$.
\end{proof}

For example, a nontrivial additive character $a(n)=e^{2\pi i kn/q}$ has
zero mean. Every normalized order-$r$ jet for this periodically weighted
spectrum, even after multiplication by $n^dP(n)$, has degree at most $r$.
There is no first-derivative multiplicative anomaly for such a pure
zero-mean periodic component. This does not say that its individual jets
vanish.

The finite parts themselves retain the periodic data. For one degree $d$,
the constant coefficient at $s=d+1$ of the corresponding term in
\eqref{asj:eq:quasipoly} is
\[
 \frac1q\sum_{j=1}^q a_d(j)\gamma_0(j/q)-\mu_d\log q.
\]
Other degrees contribute holomorphic values there. Higher coefficients
follow by multiplying $e^{-s\log q}$ by \eqref{asj:eq:stieltjes}. Thus the
mean law is a statement about the local defect, not the entire answer.

\subsection{A two-pole example from the sum-of-divisors function}
Let $\sigma_1(n)=\sum_{d\mid n}d$. The elementary product of two
absolutely convergent Dirichlet series proves
\begin{equation}\label{asj:eq:sigma1}
 \sum_{n\ge1}\frac{\sigma_1(n)}{n^s}=\zeta(s)\zeta(s-1),
 \qquad \Rea s>2.
\end{equation}
The two poles are simple, with
\[
 \kappa_{1,1}=\zeta(0)=-\frac12,
 \qquad \kappa_{2,1}=\zeta(2).
\]
The coexistence of distinct poles is genuinely different from inserting
a power of $\log n$ into one polynomially weighted spectrum.

For $P=1$, denote the continued multi-shift series by $Z_{\boldsymbol w}$,
and the single-shift series by $F_a$. Both are holomorphic at the spectral
origin, since all relevant numerators vanish there and the poles are simple.
Set
\[
 \delta_{\boldsymbol w}=
        \sum_j w_jF_{a_j}'(0)-Z_{\boldsymbol w}'(0),\qquad W\ne0.
\]
\begin{proposition}[An exact arithmetic anomaly]\label{asj:thm:sigma}
For $a_j>-1$,
\begin{equation}\label{asj:eq:sigma-anomaly}
 \boxed{\displaystyle
 \delta_{\boldsymbol w}
      =\frac{\zeta(2)}{2W}
                 \sum_{i<j}w_iw_j(a_i-a_j)^2.}
\end{equation}
\end{proposition}
\begin{proof}
Normalize the weights. For $r=1$, the nonlocal part is linear and cancels
against interpolation at the coordinate vertices. The pole at $1$ gives
$[x^{-1}]L_{\boldsymbol\alpha}^2=0$. The pole at $2$ gives
$\zeta(2)(\sum_j\alpha_ja_j)^2/2$. Therefore the difference between
vertex interpolation and the jet is
\[
 \frac{\zeta(2)}2\left\{
       \sum_j\alpha_ja_j^2-\left(\sum_j\alpha_ja_j\right)^2\right\}
   =\frac{\zeta(2)}2\sum_{i<j}\alpha_i\alpha_j(a_i-a_j)^2.
\]
Multiply by $W$ using \eqref{asj:eq:scaling}. The small-shift argument extends
to all stated shifts by the finite-head theorem.
\end{proof}

For two unit weights and shifts $a,b$, the anomaly is
$\zeta(2)(a-b)^2/4$. This is a concrete arithmetic specialization of the
classical simple-pole mechanism, not a claim to have discovered that
mechanism. In contrast, the next sections concern base functions with
higher poles, for which $F_a'(0)$ is typically not defined at all.

% END sections/04-arithmetic.tex


% BEGIN sections/05-hurwitz.tex
\section{Multiplicative Hurwitz lattices and all Stieltjes layers}
\label{asj:sec:hurwitz}

\subsection{An unconditional higher-pole family}
Fix $h\ge1$ and $b_1,\ldots,b_h>0$. Let
\begin{equation}\label{asj:eq:lattice}
 \lambda_{\boldsymbol n}=\prod_{j=1}^h(n_j+b_j),
 \qquad \boldsymbol n\in\N_0^h,
 \qquad \lambda_* =\prod_jb_j.
\end{equation}
This spectrum is discrete with finite heads. Multiplying absolutely
convergent series gives
\begin{equation}\label{asj:eq:Dhurwitz}
 \D(s)=\sum_{\boldsymbol n}\lambda_{\boldsymbol n}^{-s}
              =\prod_{j=1}^h\zeta(s,b_j),\qquad \Rea s>1.
\end{equation}
Thus there is just one pole, at $s=1$, of order exactly $h$, and its
leading coefficient is one. When all $b_j=1$, collecting equal integer
products gives
\[
 \D(s)=\sum_{N\ge1}\frac{d_h(N)}{N^s}=\zeta(s)^h,
 \qquad d_h(N)=\#\{(n_1,\ldots,n_h)\in\N^h:n_1\cdots n_h=N\}.
\]
These are nonpolynomial arithmetic multiplicities. In particular,
$d_2(N)$ is the ordinary number-of-divisors function.

Write
\begin{equation}\label{asj:eq:dcoeff}
 \D(1+s)=\sum_{v=-h}^{\infty}d_v s^v.
\end{equation}
The notation $d_v$ here is a Laurent coefficient, distinguished from the
arithmetic function $d_h(N)$ by its argument. Define
\[
 A_j(s)=1+\sum_{n\ge0}\frac{(-1)^n\gamma_n(b_j)}{n!}s^{n+1}.
\]
Then the exact finite recipe is
\begin{equation}\label{asj:eq:drecipe}
 d_v=[s^{v+h}]\prod_{j=1}^h A_j(s).
\end{equation}
For any fixed $v$ this uses only finitely many Stieltjes constants. In
particular,
\begin{align}
 d_{-h}&=1,\notag\\
 d_{-(h-1)}&=\sum_j\gamma_0(b_j),\label{asj:eq:dleading}\\
 d_{-(h-2)}&=\sum_{i<j}\gamma_0(b_i)\gamma_0(b_j)
                   -\sum_j\gamma_1(b_j).\notag
\end{align}
For two parameters $b,c$, the first regular coefficients are
\begin{align}
 d_0&=\gamma_0(b)\gamma_0(c)-\gamma_1(b)-\gamma_1(c),\label{asj:eq:d0}\\
 d_1&=\frac{\gamma_2(b)+\gamma_2(c)}2
             -\gamma_0(b)\gamma_1(c)-\gamma_1(b)\gamma_0(c),\label{asj:eq:d1}\\
 d_2&=-\frac{\gamma_3(b)+\gamma_3(c)}6
      +\frac{\gamma_0(b)\gamma_2(c)+\gamma_2(b)\gamma_0(c)}2
      +\gamma_1(b)\gamma_1(c).\label{asj:eq:d2}
\end{align}
These follow by multiplying \eqref{asj:eq:stieltjes}; no regularized product
of divergent constants is involved.

\subsection{Every local term, not only the highest pole}
For the lattice \eqref{asj:eq:lattice}, Theorem~\ref{asj:thm:transfer} becomes
\begin{equation}\label{asj:eq:Hurwitzgerm}
 \calZ_P(\boldsymbol u)
  =\calH_P(\boldsymbol u)
     +\calR_1(\boldsymbol u)\sum_{j=1}^h\frac{d_{-j}}{U^j},
 \qquad \calR_1=\sum_{k=0}^d p_k C_{k+1}.
\end{equation}
Consequently the pole coefficients are independent of the auxiliary
finite head but depend nontrivially on the Hurwitz parameters. The
coefficient $d_{-h}=1$ makes the very highest difference independent of
those parameters:
\begin{equation}\label{asj:eq:Hurwitztop}
 \Delta_{\boldsymbol v_1}\cdots\Delta_{\boldsymbol v_{r+h}}
 \calJ_r
 =(-1)^{r+h}r![x^{-1}]P(x)\prod_iV_{\boldsymbol v_i}(x).
\end{equation}
If $P$ has degree less than $r+h-1$, this particular difference vanishes.
If $P$ has degree exactly $r+h-1$ with leading coefficient $p$, it equals
\begin{equation}\label{asj:eq:Hurwitzsharp}
 (-1)^{r+h}r!p\prod_{i=1}^{r+h}\left(\sum_jv_{ij}a_j\right).
\end{equation}
Thus the degree bound is attained whenever this product is nonzero.
Lower differences generally retain the Stieltjes coefficients in
\eqref{asj:eq:dleading}; replacing the whole principal part by its leading
term would give incorrect lower-order formulas.

\subsection{A normally convergent expression for the individual jets}
Put
\[
 A_{t,\ell}(\boldsymbol\alpha)
   =[z^\ell]\left(\sum_j\alpha_j\log(1+a_jz)\right)^t,
 \quad A_{0,0}=1,\quad A_{0,\ell}=0\ (\ell>0).
\]
The following formula makes the nonlocal coordinates explicit as well.

\begin{theorem}[All individual Hurwitz-product jets]\label{asj:thm:alljets}
Suppose $\max_j|a_j|<\lambda_*$. For every $r\ge0$,
\begin{align}\label{asj:eq:alljets}
 \calJ_r(\boldsymbol\alpha)
 =\sum_{k=0}^d p_k\Bigg\{&\D^{(r)}(-k)
  +\sum_{\substack{\ell\ge1\\\ell\ne k+1}}\sum_{t=1}^r
    (-1)^t\binom rt A_{t,\ell}\D^{(r-t)}(\ell-k)\notag\\
 &+r!\sum_{t=1}^{r+h}\frac{(-1)^t}{t!}
                       A_{t,k+1}d_{r-t}\Bigg\}.
\end{align}
Every infinite sum converges normally on smaller compact shift and
exponent sets. For arbitrary allowed real shifts, the same construction
uses a finite original head and its entire subtraction from $\D$.
\end{theorem}
\begin{proof}
For the stated small shifts, take an empty head in \eqref{asj:eq:tail} and
expand
\[
 C_\ell(s\boldsymbol\alpha)
     =\sum_{t\ge0}\frac{(-1)^t}{t!}A_{t,\ell}s^t.
\]
At $\ell=0$, the contribution is $\D^{(r)}(-k)$. At nonresonant
$\ell\ge1$, multiply by the Taylor series of $\D(s+\ell-k)$.
At the unique resonant value $\ell=k+1$, multiply instead by
\eqref{asj:eq:dcoeff}. The coefficient of $s^r$ receives terms through
$t=r+h$, exactly as displayed. Lemma~\ref{asj:thm:tail} justifies every
coefficient extraction and proves the finite-head extension.
\end{proof}

The last sum is finite, but its upper endpoint is essential. Truncating
$C_{k+1}(s\boldsymbol\alpha)$ at degree $r$ before multiplying by the
pole loses all terms responsible for the local identities.

\subsection{A direct obstruction to using ordinary derivatives}
For one shifted factor and $P=1$, write
\[
 F_a(s)=\sum_{\boldsymbol n}(\lambda_{\boldsymbol n}+a)^{-s}.
\]
Since $C_1(s)=-as$, the only polar contribution at zero is
\begin{equation}\label{asj:eq:Fpolar}
 \operatorname{pp}_{s=0}F_a(s)
       =-a\sum_{j=2}^h d_{-j}s^{1-j}.
\end{equation}
For $h=2$, in particular,
\begin{equation}\label{asj:eq:doublepole-example}
 \Res_{s=0}F_a(s)=-a,\qquad
 [s^0]F_a(s)=\D(0)-a\{\gamma_0(b)+\gamma_0(c)\}.
\end{equation}
Thus $F_a'(0)$ does not exist for $a\ne0$. This is a proved warning
against an overextension of the ordinary determinant convention, not
an error attributed to the earlier report, which already distinguished
Laurent coefficients for logarithmic multiplicities.

% END sections/05-hurwitz.tex


% BEGIN sections/06-alternants.tex
\section{Exact alternating identities for arithmetic Laurent jets}
\label{asj:sec:alternants}

Let $m\ge2$,
\[
 N=\binom m2,\qquad K_m=\prod_{k=0}^{m-1}k!,\qquad
 \Delta(\boldsymbol x)=\prod_{i<j}(x_j-x_i).
\]
We use $(\sigma\boldsymbol w)_j=w_{\sigma(j)}$ and
$W=\sum_jw_j\ne0$. For the multiplicative Hurwitz lattice of order $h$
assume $h\le N$, and set $r=N-h$.

\begin{theorem}[Higher-pole spectral alternant]\label{asj:thm:alternant}
For every polynomial $P$ and allowed shifts,
\begin{align}\label{asj:eq:alternant}
 &\sum_{\sigma\in S_m}\operatorname{sgn}(\sigma)
           r![s^r]\calZ_P(s\sigma\boldsymbol w)\notag\\
 &\quad=\frac{(-1)^N r!}{W^h K_m}\Delta(\boldsymbol w)
          [x^{-1}]P(x)\prod_{i<j}\log\frac{x+a_j}{x+a_i}.
\end{align}
If $P$ has degree $N-1$ and leading coefficient $p$, this reduces to
\begin{equation}\label{asj:eq:alternant-leading}
 \boxed{\displaystyle
 \sum_{\sigma}\operatorname{sgn}(\sigma)
           r![s^r]\calZ_P(s\sigma\boldsymbol w)
  =\frac{(-1)^N r!p}{W^hK_m}
                \Delta(\boldsymbol w)\Delta(\boldsymbol a).}
\end{equation}
For degree below $N-1$, the alternating sum is zero.
\end{theorem}
\begin{proof}
An alternating polynomial in $m$ variables is divisible by their
Vandermonde, whose degree is $N$. In \eqref{asj:eq:Jlocal}, alternation
therefore kills the holomorphic term of degree $r<N$ and every local
term of degree $r+j<N$. Only the top principal coefficient $d_{-h}=1$
remains. The elementary identity
\begin{equation}\label{asj:eq:elementary-alternant}
 \sum_{\sigma}\operatorname{sgn}(\sigma)
       \left(\sum_j\alpha_{\sigma(j)}y_j\right)^N
       =\frac{N!}{K_m}\Delta(\boldsymbol\alpha)\Delta(\boldsymbol y)
\end{equation}
follows by expanding the power. Terms with a repeated exponent vanish;
the only distinct nonnegative exponents summing to $N$ are
$0,1,\ldots,m-1$. Their two determinants give the Vandermondes with the
orientation displayed. Substituting $y_j=\log(1+a_j/x)$ in
\eqref{asj:eq:Rhom} proves the normalized identity. Equation
\eqref{asj:eq:scaling} and $\Delta(\boldsymbol w/W)=W^{-N}\Delta(\boldsymbol w)$
produce $W^{r-N}=W^{-h}$. Finally,
\[
 \prod_{i<j}\log\frac{x+a_j}{x+a_i}
         =\Delta(\boldsymbol a)x^{-N}+O(x^{-N-1}),
\]
which gives the last two assertions.
\end{proof}

\subsection{A six-term first-jet identity for the divisor function}
For a monic quadratic $P$ define
\begin{equation}\label{asj:eq:Tdef}
 T_P(a,b)=[s]\sum_{n\ge1}d_2(n)P(n)
                  \bigl((n+a)^2(n+b)\bigr)^{-s},\qquad a,b>-1,
\end{equation}
by meromorphic continuation and the fixed Laurent convention.

\begin{corollary}[Six first jets collapse to an algebraic number]
\label{asj:thm:six}
\begin{align}\label{asj:eq:six}
 &T_P(a,b)+T_P(b,c)+T_P(c,a)\notag\\
 &\quad-T_P(b,a)-T_P(c,b)-T_P(a,c)
       =\frac{(a-b)(b-c)(c-a)}9.
\end{align}
The lower two coefficients of $P$ do not affect the answer. For a
quadratic leading coefficient $p$, multiply the right side by $p$;
for degree at most one the answer is zero.
\end{corollary}
\begin{proof}
Apply Theorem~\ref{asj:thm:alternant} with $m=3$, $h=2$, $r=1$,
$\boldsymbol w=(2,1,0)$, $W=3$, $K_3=2$, and
$\Delta(\boldsymbol w)=-2$. Reading the six permutations in
\eqref{asj:eq:alternant-leading} gives precisely \eqref{asj:eq:six}.
\end{proof}

\paragraph{A compact direct certificate.}
For the six permutations of $(2,1,0)$ let $\mathcal A$ denote signed
summation. Since $C_\ell(s\boldsymbol w)$ is polynomial in the weights,
its alternating part is divisible by $s^3$. It vanishes identically for
$\ell<3$. At $\ell=3$, equation~\eqref{asj:eq:elementary-alternant} gives
\[
 \mathcal A C_3(s\boldsymbol w)=s^3\Delta(a,b,c).
\]
For $P(x)=x^2$, only $\ell=3$ can contribute to the first Laurent
coefficient: all other nonzero alternating terms multiply a function
holomorphic at $s=0$. The resonant factor is
$\zeta(1+3s)^2=1/(9s^2)+O(s^{-1})$. Hence its first coefficient is
$\Delta(a,b,c)/9$. The lower coefficients of $P$ resonate only at
$\ell<3$ and make no contribution. This is a finite symbolic certificate
of the corollary independent of any numerical spectral evaluation.

For $(a,b,c)=(1,2,3)$ the value is $2/9$. Replacing $d_2(n)$ by the
three-factor divisor multiplicity $d_3(n)$ moves the same identity down
one more spectral order: with
\[
 U_P(a,b)=[s^0]\sum_{n\ge1}d_3(n)P(n)
                  \bigl((n+a)^2(n+b)\bigr)^{-s},
\]
the same signed six-term combination equals
\begin{equation}\label{asj:eq:sixconstant}
 \frac{(a-b)(b-c)(c-a)}{27}.
\end{equation}
It is an identity between constant Laurent coefficients, not between
values of functions at their poles.

\subsection{A hierarchy with independent Hurwitz parameters}
The identities \eqref{asj:eq:six} and \eqref{asj:eq:sixconstant} hold unchanged
when the integer divisor spectrum is replaced by arbitrary multiplicative
Hurwitz lattices of two and three factors, respectively. Although each
individual Laurent jet depends on all the $b_j$, the alternating
combination does not. The reason is exactly $d_{-h}=1$, not a cancellation
suggested by numerical fitting.

For $m=4$, $N=6$ and $K_4=12$. With $h=2$, a signed sum of twenty-four
fourth jets for a degree-five polynomial with leading coefficient $p$
is
\[
 \frac{2p}{W^2}\Delta(\boldsymbol w)\Delta(\boldsymbol a).
\]
With $h=3$, the corresponding third-jet sum is
$p\Delta(\boldsymbol w)\Delta(\boldsymbol a)/(2W^3)$.
This makes visible the general shift in spectral order: the same
Vandermonde identity occurs at order $N-h$ when the base pole has order
$h$. These are exact consequences of the higher-pole germ; applying a
simple-pole determinant formula at an undefined derivative would not
establish them.

% END sections/06-alternants.tex


% BEGIN sections/07-gamma.tex
\section{Gamma reciprocity, polygamma sums, and a normalized primitive}
\label{asj:sec:gamma}

\subsection{A convergent nonlinear part of the first Laurent jet}
For the multiplicative Hurwitz lattice define
\begin{equation}\label{asj:eq:Edef}
 E_{\boldsymbol b}(a)=\sum_{\boldsymbol n}
        \left\{\frac a{\lambda_{\boldsymbol n}}
                    -\log\left(1+\frac a{\lambda_{\boldsymbol n}}\right)\right\},
 \qquad a> -\lambda_*.
\end{equation}
Every brace must be kept grouped. Its tail is $O(\lambda_{\boldsymbol n}^{-2})$
on compact shift sets, and $\sum\lambda_{\boldsymbol n}^{-2}
=\prod_j\zeta(2,b_j)<\infty$. Thus this is an ordinary, locally normally
convergent series, not a finite-part sum.

\begin{proposition}[First-jet decomposition]\label{asj:thm:firstjet}
For $F_a$ of Section~\ref{asj:sec:hurwitz},
\begin{equation}\label{asj:eq:firstjet}
 \boxed{\displaystyle \calJ_1F_a=\D'(0)-a d_0+E_{\boldsymbol b}(a).}
\end{equation}
\end{proposition}
\begin{proof}
For $|a|<\lambda_*$, apply \eqref{asj:eq:alljets} with $P=1$, one exponent,
and $r=1$. The resonant term is $-a d_0$. For $\ell\ge2$,
$[s](s)_\ell/\ell!=1/\ell$, so the remaining sum is
\[
 \sum_{\ell\ge2}\frac{(-a)^\ell}{\ell}\D(\ell)
       =\sum_{\boldsymbol n}\left\{a/\lambda_{\boldsymbol n}
                  -\log(1+a/\lambda_{\boldsymbol n})\right\}.
\]
Absolute convergence justifies the interchange. Both sides continue
locally holomorphically throughout the slit shift plane
$\C\setminus(-\infty,-\lambda_*]$, using finite heads on the spectral
side and the grouped series on the other side. The identity theorem,
or continuation along the real interval, gives every stated $a$.
\end{proof}

The entire principal part at the spectral origin is affine in $a$ by
\eqref{asj:eq:Fpolar}. All nonlinearity of the first Laurent coefficient is
therefore the convergent function $E_{\boldsymbol b}$.

\subsection{Two independent Gamma representations of the same function}
For two parameters write $E_{b,c}=E_{(b,c)}$.

\begin{theorem}[Centered Gamma reciprocity]\label{asj:thm:gamma}
For $b,c>0$ and $a>-bc$,
\begin{align}\label{asj:eq:gamma-reciprocity}
 E_{b,c}(a)
 &=\sum_{n=0}^{\infty}\left\{
       \log\Gamma\left(c+\frac a{n+b}\right)-\log\Gamma(c)
                      -\frac{a\psi(c)}{n+b}\right\}\notag\\
 &=\sum_{m=0}^{\infty}\left\{
       \log\Gamma\left(b+\frac a{m+c}\right)-\log\Gamma(b)
                      -\frac{a\psi(b)}{m+c}\right\}.
\end{align}
Both series converge absolutely and locally uniformly in $a$. On
$|a|<bc$ they also equal
\begin{equation}\label{asj:eq:gamma-zeta}
 E_{b,c}(a)=\sum_{\ell\ge2}
                 \frac{(-a)^\ell}{\ell}\zeta(\ell,b)\zeta(\ell,c).
\end{equation}
\end{theorem}
\begin{proof}
The classical centered Gamma product gives, for $c>0$ and $t>-c$,
\begin{equation}\label{asj:eq:centeredgamma}
 \log\Gamma(c+t)-\log\Gamma(c)-t\psi(c)
    =\sum_{m\ge0}\left\{\frac{t}{m+c}
                         -\log\left(1+\frac t{m+c}\right)\right\}.
\end{equation}
One may derive it without choosing a product normalization: both sides
vanish with their first derivative at $t=0$, and their second derivatives
are the same convergent trigamma series. Set $t=a/(n+b)$ and sum over
$n$. All double summands are absolutely summable, since their tails are
bounded by a constant times $(n+b)^{-2}(m+c)^{-2}$ and their finitely many
small factors stay away from zero on compact allowed shift sets.
Fubini's theorem gives $E_{b,c}$. Interchanging $m$ and $n$ proves the
second representation. Expanding the logarithm under $|a|<bc$ proves
\eqref{asj:eq:gamma-zeta}.
\end{proof}

For $b=c=1$, the first Laurent coefficient becomes the explicit
Gamma--Stieltjes identity
\begin{equation}\label{asj:eq:divisor-gamma}
 \calJ_1F_a=\frac12\log(2\pi)-a(\gamma^2-2\gamma_1)
      +\sum_{n\ge1}\left\{\log\Gamma(1+a/n)+\frac{\gamma a}{n}\right\}.
\end{equation}
Here $\D'(0)=2\zeta(0)\zeta'(0)=\tfrac12\log(2\pi)$ and
$d_0=\gamma^2-2\gamma_1$. At the same time,
\[
 \Res_{s=0}F_a=-a,\qquad [s^0]F_a=\frac14-2\gamma a.
\]
The Gamma series in \eqref{asj:eq:divisor-gamma} cannot make that nonzero pole
 disappear. It evaluates the correctly specified Laurent coefficient.

\subsection{Reciprocal digamma and all polygamma derivatives}
\begin{corollary}[All argument derivatives]\label{asj:thm:polygamma}
For $b,c>0$, $a>-bc$,
\begin{align}\label{asj:eq:digamma-reciprocity}
 E_{b,c}'(a)
 &=\sum_{n\ge0}\frac{\psi(c+a/(n+b))-\psi(c)}{n+b}\notag\\
 &=\sum_{m\ge0}\frac{\psi(b+a/(m+c))-\psi(b)}{m+c}.
\end{align}
For every integer $r\ge2$,
\begin{align}\label{asj:eq:polygamma-reciprocity}
 E_{b,c}^{(r)}(a)
 &=\sum_{n\ge0}\frac{\psi^{(r-1)}(c+a/(n+b))}{(n+b)^r}\notag\\
 &=\sum_{m\ge0}\frac{\psi^{(r-1)}(b+a/(m+c))}{(m+c)^r}\notag\\
 &=(-1)^r(r-1)!\sum_{n,m\ge0}
                       \bigl((n+b)(m+c)+a\bigr)^{-r}.
\end{align}
At $a=0$ this is $(-1)^r(r-1)!\zeta(r,b)\zeta(r,c)$.
\end{corollary}
\begin{proof}
Differentiate the grouped normally convergent series on compact subsets
of the allowed shift domain. The $r$th derivative for $r\ge2$ of one
grouped summand in \eqref{asj:eq:Edef} is
$(-1)^r(r-1)!(\lambda+a)^{-r}$. The Gamma representation gives the
first two forms. The conventional polygamma series, or differentiation
of \eqref{asj:eq:centeredgamma}, gives the last one. At $a=0$ absolute
convergence factors the double series.
\end{proof}

For more than two Hurwitz factors, summing over one factor gives an
analogous centered Gamma series over the remaining $(h-1)$ indices.
The choice of the factor summed first is arbitrary, producing $h$
compatible Gamma representations. This is a constructive representation,
not a reduction claim into finitely many Gamma values.

\subsection{An anchored primitive in Hurwitz-zeta derivatives}
Define the classical primitive
\begin{equation}\label{asj:eq:Kdef}
 K(z)=\zeta'(-1,z)+\frac{z-z^2}{2}+\frac z2\log(2\pi),\qquad z>0.
\end{equation}
The identities $\partial_z\zeta(s,z)=-s\zeta(s+1,z)$ and
$\zeta'(0,z)=\log\Gamma(z)-\tfrac12\log(2\pi)$ imply
$K'(z)=\log\Gamma(z)$ \cite{dlmf-hurwitz}. This agrees with the
primitive already used in the manuscript \cite{repo}.

\begin{theorem}[Reciprocal normalized primitive]\label{asj:thm:primitive}
Let $I_{b,c}(a)=\int_0^a E_{b,c}(t)\dd t$, with oriented integration
when $a<0$. Then
\begin{align}\label{asj:eq:primitive-gamma}
 I_{b,c}(a)=\sum_{n\ge0}\Bigg\{&(n+b)
              \left[K\left(c+\frac a{n+b}\right)-K(c)\right]\notag\\
              &-a\log\Gamma(c)-\frac{a^2\psi(c)}{2(n+b)}\Bigg\}.
\end{align}
The same expression with $b$ and $c$ exchanged is equal to it. Equivalently,
\begin{align}\label{asj:eq:primitive-lattice}
 I_{b,c}(a)=\sum_{n,m\ge0}\Bigg\{&\frac{a^2}{2\lambda_{n,m}}
        -(\lambda_{n,m}+a)
                 \log\left(1+\frac a{\lambda_{n,m}}\right)+a\Bigg\},
\end{align}
where $\lambda_{n,m}=(n+b)(m+c)$. For $|a|<bc$,
\begin{equation}\label{asj:eq:primitive-power}
 I_{b,c}(a)=\sum_{\ell\ge2}\frac{(-1)^\ell a^{\ell+1}}{\ell(\ell+1)}
                                   \zeta(\ell,b)\zeta(\ell,c).
\end{equation}
\end{theorem}
\begin{proof}
Integrate the grouped summands of Theorem~\ref{asj:thm:gamma} along the
compact segment from $0$ to $a$. Uniform convergence permits the
interchange. Substitution $u=c+t/(n+b)$ gives the $K$ term, and the
two subtractions integrate to the other two terms in
\eqref{asj:eq:primitive-gamma}. Integration of \eqref{asj:eq:Edef} gives
\eqref{asj:eq:primitive-lattice}. Its grouped tail is $O(\lambda_{n,m}^{-2})$.
Both forms vanish at $a=0$, fixing the integration constant. Integrating
\eqref{asj:eq:gamma-zeta} proves \eqref{asj:eq:primitive-power}.
\end{proof}

In \eqref{asj:eq:primitive-gamma}, the three large pieces must not be summed
separately. The cancellations are part of the identity's definition.
This is the same normalization issue that occurs in negative polygamma,
but here both the base point and every convergent subtraction are explicit.

% END sections/07-gamma.tex


% BEGIN sections/08-harmonic.tex
\section{Generalized harmonic numbers in every single-shift jet}
\label{asj:sec:harmonic}

Let $e_j(q)$ be the elementary symmetric polynomial of degree $j$ in
$1,1/2,\ldots,1/q$, with $e_0(q)=1$ and $e_j(q)=0$ for $j>q$.
For $q=0$ use the empty alphabet. Define
$H_q^{(j)}=\sum_{n=1}^q n^{-j}$ and $H_q=H_q^{(1)}$.
The finite product gives
\begin{align}\label{asj:eq:harmonic-egf}
 \frac{(s)_\ell}{\ell!}
 &=\frac{s}{\ell}\prod_{j=1}^{\ell-1}(1+s/j)
   =\frac{s}{\ell}\sum_{j=0}^{\ell-1}e_j(\ell-1)s^j,\notag\\
 \sum_{j\ge0}e_j(q)s^j
 &=\exp\left(\sum_{k\ge1}\frac{(-1)^{k+1}}kH_q^{(k)}s^k\right).
\end{align}
The second equality is understood as a power-series identity near zero;
its exponential simplifies to the finite product. In particular,
\[
 e_1(q)=H_q,\qquad
 e_2(q)=\tfrac12\{H_q^2-H_q^{(2)}\},\qquad
 e_3(q)=\tfrac16\{H_q^3-3H_qH_q^{(2)}+2H_q^{(3)}\}.
\]
These are the familiar harmonic/Stirling coordinates already present in
the manuscript's derivative tower. Their role here is to control
arithmetic spectral finite parts, not just parameter derivatives of a
single Hurwitz function.

\begin{theorem}[All single-shift Laurent jets]\label{asj:thm:harmonic}
For the multiplicative Hurwitz lattice, $|a|<\lambda_*$, and $r\ge1$,
\begin{align}\label{asj:eq:single-all}
 \calJ_r F_a
  ={}&\D^{(r)}(0)-a r!d_{r-1}\notag\\
   &+\sum_{\ell\ge2}\frac{(-a)^\ell}{\ell}
       \sum_{t=1}^{\min(r,\ell)}
         \frac{r!}{(r-t)!}\,e_{t-1}(\ell-1)
                                  \D^{(r-t)}(\ell).
\end{align}
The constant Laurent coefficient is
$\calJ_0F_a=\D(0)-a d_{-1}$. All spectral derivatives of $\D$
are finite sums of products of Hurwitz-zeta derivatives by Leibniz's
rule. Thus \eqref{asj:eq:single-all} uses only specified zeta jets,
generalized harmonic numbers, and finite Stieltjes polynomials.
\end{theorem}
\begin{proof}
The binomial expansion is
\[
 F_a(s)=\sum_{\ell\ge0}\frac{(-a)^\ell(s)_\ell}{\ell!}\D(s+\ell).
\]
The $\ell=0$ term contributes $\D^{(r)}(0)$. The resonant term $\ell=1$
is $-as\D(1+s)$ and contributes $-a r!d_{r-1}$, including $r=0$.
For $\ell\ge2$, substitute \eqref{asj:eq:harmonic-egf} and multiply by
the ordinary Taylor series of $\D(\ell+s)$. Theorem~\ref{asj:thm:alljets}
justifies the extraction and convergence.
\end{proof}

The first genuinely higher rows are worth displaying without compressed
symmetric-function notation:
\begin{align}\label{asj:eq:secondjet}
 \calJ_2F_a={}&\D''(0)-2a d_1
     +2\sum_{\ell\ge2}\frac{(-a)^\ell}{\ell}
           \{\D'(\ell)+H_{\ell-1}\D(\ell)\},\\
 \calJ_3F_a={}&\D'''(0)-6a d_2
     +3\sum_{\ell\ge2}\frac{(-a)^\ell}{\ell}
        \Big\{\D''(\ell)+2H_{\ell-1}\D'(\ell)\notag\\
        &\hspace{43mm}
        +(H_{\ell-1}^2-H_{\ell-1}^{(2)})\D(\ell)\Big\}.
 \label{asj:eq:thirdjet}
\end{align}
For two independent Hurwitz parameters, $d_1,d_2$ are exactly
\eqref{asj:eq:d1}--\eqref{asj:eq:d2}. For example,
\[
 \D'(\ell)=\zeta'(\ell,b)\zeta(\ell,c)
                    +\zeta(\ell,b)\zeta'(\ell,c),
\]
so every occurrence in \eqref{asj:eq:secondjet} is explicit. These formulas
are exact convergent identities; no assertion is made that the infinite
series always collapses to finitely many familiar constants.

\subsection{A finite algorithm and its shift continuation}
To compute a jet of order $r$, one first computes the finite Laurent
coefficients through $d_{r-1}$ using \eqref{asj:eq:drecipe}, then generates
the elementary symmetric coefficients through degree $r-1$, and finally
evaluates the normally convergent zeta-jet series. Only finitely many
algebraic operations are needed at each series index. The recurrence
\begin{equation}\label{asj:eq:Crecurrence}
 \ell C_\ell(\boldsymbol u)
     =\sum_{j=1}^{\ell}(-1)^j
         \left(\sum_i u_i a_i^j\right)C_{\ell-j}(\boldsymbol u)
\end{equation}
follows by differentiating \eqref{asj:eq:Cdef} with respect to $z$ and
is convenient for multiple factors. The verification code tests it
against the independent binomial product \eqref{asj:eq:Cpoch}.

The radius $|a|<\lambda_*$ is a convergence condition for this particular
Taylor series, not a singularity of the first-jet Gamma representation at
positive $a$. For larger shifts use the finite-head version of the transfer
theorem, or use \eqref{asj:eq:firstjet} and its derivatives. Confusing a
convergent expansion's radius with the full identity's domain would
unnecessarily discard most positive shifts.

% END sections/08-harmonic.tex


% BEGIN sections/09-twists.tex
\section{Polylogarithmic twists and exact finite-part untwisting}
\label{asj:sec:twists}

\subsection{Twisting lowers the pole order}
For $0<z<1$ and $c>0$, the defining Lerch series
\[
 \Phi(z,s,c)=\sum_{n\ge0}\frac{z^n}{(n+c)^s}
\]
is entire in $s$, by uniform exponential convergence on compact sets
\cite{dlmf-lerch}. Consequently the base series of a multiplicative lattice
with $H$ untwisted factors and some geometrically twisted factors is
\begin{equation}\label{asj:eq:twistedbase}
 \D_{\boldsymbol z}(s)=
      \prod_{j=1}^H\zeta(s,b_j)\prod_{k=1}^v\Phi(z_k,s,c_k).
\end{equation}
For real $0<z_k<1$, its pole order at $1$ is exactly $H$ and its leading
coefficient is
\[
 \kappa_{1,H}=\prod_{k=1}^v\Phi(z_k,1,c_k)>0.
\]
With $H=0$ it is entire. Complex twists are also covered by the transfer
theorem, but a zero of a Lerch factor can lower the pole order further;
one must then use the actual principal part, not the number of nominally
untwisted factors.

For a particularly transparent arithmetic case, put
\begin{equation}\label{asj:eq:cz}
 c_z(N)=\sum_{d\mid N}z^d d_2(N/d),\qquad
 \D_z(s)=\sum_{N\ge1}\frac{c_z(N)}{N^s}
                         =\zeta(s)^2\Li_s(z).
\end{equation}
The equality follows by an absolutely convergent Dirichlet convolution.
For $0<z<1$, the leading double-pole coefficient is
$\Li_1(z)=-\log(1-z)$. Define $T_{P,z}$ as in \eqref{asj:eq:Tdef}, with
$d_2(N)$ replaced by $c_z(N)$. The proof of Theorem~\ref{asj:thm:alternant}
gives the exact identity
\begin{align}\label{asj:eq:twisted-six}
 &T_{P,z}(a,b)+T_{P,z}(b,c)+T_{P,z}(c,a)\notag\\
 &\quad-T_{P,z}(b,a)-T_{P,z}(c,b)-T_{P,z}(a,c)
    =-\frac{\log(1-z)}9(a-b)(b-c)(c-a)
\end{align}
for monic quadratic $P$. This directly connects arithmetic Laurent jets
to a classical polylogarithm, while all lower residue coefficients cancel.

\subsection{A concrete failure of naive untwisting}
At $z=1$, $c_z(N)$ becomes $d_3(N)$ and the base function is $\zeta(s)^3$.
Its order-three pole has subleading coefficient $3\gamma$. For the same
\emph{first} Laurent jets and the same quadratic $P$, alternation kills
the degree-one and degree-two parts. The degree-four part is zero because
$\calR_1$ has degree at most three. The remaining degree-three part gives
\begin{equation}\label{asj:eq:untwisted-first-six}
 \operatorname{Alt}T_{P,1}(a,b,c)
           =\frac\gamma3(a-b)(b-c)(c-a).
\end{equation}
Here $\operatorname{Alt}$ is the signed six-term expression in
\eqref{asj:eq:twisted-six}. Equations \eqref{asj:eq:twisted-six} and
\eqref{asj:eq:untwisted-first-six} exhibit noncommutation without any
numerical approximation: the naive $z\uparrow1$ limit diverges whenever
the three shifts are distinct, whereas the untwisted Laurent coefficient
is finite. The next theorem supplies the missing compensation at all
orders, not only in this example.

\subsection{A finite compensation for every polynomial and every jet}
Let
\[
 B(s)=\prod_{j=1}^H\zeta(s,b_j),\qquad H\ge1,
 \quad Y_{\boldsymbol n}=\prod_{j=1}^H(n_j+b_j),
 \quad \lambda_{n,\boldsymbol n}=nY_{\boldsymbol n},\quad n\ge1.
\]
The untwisted spectrum has base $\D_0(s)=\zeta(s)B(s)$ and minimum
$\lambda_* =\prod_jb_j$. For $t>0$, attach the weight $e^{-tn}$ to its
first factor. The twisted base is
\[
 \D_t(s)=\Li_s(e^{-t})B(s).
\]
For a polynomial $P(x)=\sum_{k=0}^d p_kx^k$, allowed real shifts
$a_j>-\lambda_*$, and normalized exponents $\sum\alpha_j=1$, define
\begin{equation}\label{asj:eq:Zt}
 Z_{t,P,\boldsymbol\alpha}(s)
   =\sum_{n\ge1,\boldsymbol n\ge0}
       e^{-tn}P(\lambda_{n,\boldsymbol n})
         \prod_j(\lambda_{n,\boldsymbol n}+a_j)^{-s\alpha_j}
\end{equation}
initially in its absolute-convergence half-plane and subsequently by
continuation. At $t=0$ omit the exponential weight. Use $C_\ell$ of
\eqref{asj:eq:Cdef} for these shifts.

\begin{theorem}[All-order compensated Abel untwisting]\label{asj:thm:untwist}
Define the finite meromorphic compensation
\begin{align}\label{asj:eq:general-counter}
 M_{t,P,\boldsymbol\alpha}(s)
  =\sum_{k=0}^d p_k\sum_{\ell=0}^{k+1}
      C_\ell(s\boldsymbol\alpha)
      \Gamma(1-s-\ell+k)\,
       t^{s+\ell-k-1} B(s+\ell-k).
\end{align}
There is a sufficiently small fixed circle $|s|=\rho>0$ on which
\begin{equation}\label{asj:eq:uniform-untwist}
 Z_{t,P,\boldsymbol\alpha}(s)-M_{t,P,\boldsymbol\alpha}(s)
      \longrightarrow Z_{0,P,\boldsymbol\alpha}(s)
       \quad(t\downarrow0)
\end{equation}
uniformly. In particular, for every integer $r\ge0$,
\begin{equation}\label{asj:eq:jet-untwist}
 \boxed{\displaystyle
 \calJ_r Z_{0,P,\boldsymbol\alpha}
   =\lim_{t\downarrow0}\left\{
      \calJ_r Z_{t,P,\boldsymbol\alpha}
                 -r![s^r]M_{t,P,\boldsymbol\alpha}(s)\right\}.}
\end{equation}
The number of compensation terms depends on the degree of $P$, not on
the requested jet order. For each fixed order its Laurent coefficients
are finite expressions in Gamma jets, powers of $\log t$, zeta jets,
and generalized Stieltjes constants.
\end{theorem}
\begin{proof}
First suppose $\max|a_j|<\lambda_*$. The binomial expansion gives
\[
 Z_{t,P,\boldsymbol\alpha}(s)
  =\sum_{k=0}^d p_k\sum_{\ell\ge0}C_\ell(s\boldsymbol\alpha)
          \Li_{s+\ell-k}(e^{-t}) B(s+\ell-k).
\]
Choose $0<\rho<1/4$ small enough to avoid any other local poles; the
circle never passes through an integer argument of the Gamma factors.
The classical polylogarithm expansion is \cite[Eq.~25.12.12]{dlmf-polylog}
\begin{equation}\label{asj:eq:Jonquiere}
 \Li_v(e^{-t})=\Gamma(1-v)t^{v-1}
                 +\sum_{j\ge0}\zeta(v-j)\frac{(-t)^j}{j!},
 \qquad 0<t<2\pi,
\end{equation}
with colliding terms combined by continuation when $v$ is an integer.
For each of the finitely many indices $\ell\le k+1$, the difference
between the left side and its Gamma term tends uniformly on our circle
to $\zeta(s+\ell-k)$. Multiplying by the fixed meromorphic factors,
which are bounded on the circle, gives exactly the finite subtraction
\eqref{asj:eq:general-counter}.

For $\ell\ge k+2$, no subtraction is needed. The real part of
$s+\ell-k$ exceeds $1$, and the absolutely convergent series gives
\[
 \Li_{s+\ell-k}(e^{-t})B(s+\ell-k)
    \longrightarrow \zeta(s+\ell-k)B(s+\ell-k)
\]
uniformly on the circle, by dominated convergence. The high-$\ell$
bounds from Lemma~\ref{asj:thm:tail} apply uniformly in $t$, because
$0<e^{-tn}\le1$. Together with the Cauchy bound for $C_\ell$, they give
a common summable geometric majorant. Therefore the entire tail tends
uniformly to its untwisted counterpart. This proves
\eqref{asj:eq:uniform-untwist} for small shifts.

For general allowed shifts, remove a finite spectral head so that its
tail minimum exceeds $\max|a_j|$. Each subtracted head term is entire
in $s$ and converges uniformly to the corresponding untwisted head term
as $t\downarrow0$. Apply the same argument to the tail. The finite
subtractions in the base functions converge without a compensation and
do not alter the Gamma terms in \eqref{asj:eq:general-counter}. This proves
the general case.

Finally, the Laurent coefficient is the Cauchy integral
\[
 [s^r]F(s)=\frac1{2\pi i}\int_{|s|=\rho}\frac{F(s)}{s^{r+1}}\dd s.
\]
Uniform convergence on this fixed circle proves
\eqref{asj:eq:jet-untwist}. Expanding the finitely many Gamma and $B$ factors
at their integer arguments gives the stated coefficient class. At a
pole of $B$, its finite required Laurent coefficients are precisely
finite Stieltjes polynomials.
\end{proof}

The theorem does not interchange a divergent pointwise limit at $s=0$
with differentiation. It takes a uniform limit on a contour where the
functions are defined, after an explicit finite subtraction. This
specifies both the regulator and the correction.

\subsection{A closed Bell-polynomial compensation for one shifted factor}
For $P=1$, one shifted factor, and exponent $s$, the two terms in
\eqref{asj:eq:general-counter} simplify to
\begin{equation}\label{asj:eq:single-counter}
 M_{t,a}(s)=\Gamma(1-s)t^{s-1}B(s)
                 +a\Gamma(1-s)t^sB(1+s).
\end{equation}
The sign of the second term follows from
$(-as)\Gamma(-s)=a\Gamma(1-s)$.
Write
\[
 B(1+s)=\sum_{q=-H}^{\infty}b_qs^q,\qquad
 g_j(t)=\left.\frac{\dd^j}{\dd s^j}\{\Gamma(1-s)t^s\}\right|_{s=0}.
\]
The Gamma expansion \cite{dlmf-gamma} gives
\begin{equation}\label{asj:eq:gBell}
 \Gamma(1-s)t^s
   =\exp\left((\gamma+\log t)s
                     +\sum_{j\ge2}\frac{\zeta(j)}j s^j\right).
\end{equation}
Thus $g_j$ is the complete exponential Bell polynomial in
$\gamma+\log t,1!\zeta(2),2!\zeta(3),\ldots$; this generating identity
also defines it without Bell-polynomial notation.

\begin{corollary}[Explicit compensation at arbitrary order]
\label{asj:thm:bellcounter}
For $r\ge0$,
\begin{align}\label{asj:eq:Bellcounter}
 r![s^r]M_{t,a}(s)
   ={}&\frac1t\sum_{j=0}^r\binom rj g_j(t)B^{(r-j)}(0)\notag\\
     &+a r!\sum_{q=-H}^{r}b_q\frac{g_{r-q}(t)}{(r-q)!}.
\end{align}
\end{corollary}
\begin{proof}
Multiply the Taylor series of \eqref{asj:eq:gBell} by the ordinary Taylor
series of $B(s)$ in the first term of \eqref{asj:eq:single-counter}, and
by its Laurent series $B(1+s)$ in the second. The latter product
requires Gamma coefficients through order $r+H$.
\end{proof}

For $B(s)=\zeta(s,c)$, $H=1$, and $r=1$, put
$\ell_t=\gamma+\log t$. The complete compensation is
\begin{equation}\label{asj:eq:first-counter}
 \frac{\zeta'(0,c)+\ell_t\zeta(0,c)}t
     +a\left\{\frac{\ell_t^2+\zeta(2)}2
                    +\gamma_0(c)\ell_t-\gamma_1(c)\right\}.
\end{equation}
Subtracting this from $\calJ_1Z_{t,1}$ and taking $t\downarrow0$ yields
exactly the Gamma--Stieltjes value \eqref{asj:eq:firstjet} with parameters
$(1,c)$. The $a\ell_t^2/2$ term is indispensable; it comes from the
pole of $\zeta(1+s,c)$, and a first-order truncation of the Gamma factor
would discard it.

\subsection{An explicit compensation check on the six-term family}
Return to \eqref{asj:eq:cz}, set $z=e^{-t}$, and write
$V=(a-b)(b-c)(c-a)$. Alternating the first Laurent coefficient of the
finite compensation \eqref{asj:eq:general-counter}, with $B=\zeta^2$ and
$P$ monic quadratic, gives
\begin{equation}\label{asj:eq:altcounter}
 \operatorname{Alt}[s]M=-\frac{\log t+3\gamma}{9}\,V.
\end{equation}
To check the sign and constant directly, use unnormalized weights
$(2,1,0)$. The alternated coefficient $C_3(s\boldsymbol w)$ is
$s^3V$, while all $C_\ell$ with $\ell<3$ have zero alternation.
The only relevant compensation is
\[
 s^3V\,\Gamma(-3s)t^{3s}\zeta(1+3s)^2.
\]
Its first coefficient is \eqref{asj:eq:altcounter}. Hence the compensated
signed first jets satisfy the exact finite-$t$ identity
\begin{equation}\label{asj:eq:altfinite}
 \operatorname{Alt}T_{P,e^{-t}}+\frac{\log t+3\gamma}{9}V
  =\frac\gamma3V+\frac{V}{9}\log\frac{t}{1-e^{-t}}.
\end{equation}
The last logarithm tends to zero. This elementary special case verifies
both the divergent subtraction and the finite Euler-constant correction
in the general contour argument.

% END sections/09-twists.tex


% BEGIN sections/10-audit.tex
\section{Source disposition, verification, and further research}
\label{asj:sec:audit}

\subsection{What has been completed}
The question Q10 in \cite{prior} is resolved for the concrete classes
requested there: periodic and quasi-polynomial multiplicities are
handled by \eqref{asj:eq:quasipoly} and Theorem~\ref{asj:thm:mean}; several
poles and higher poles are handled by Theorems~\ref{asj:thm:transfer}--\ref{asj:thm:all-differences}
under explicit continuation hypotheses, with unconditional arithmetic
examples \eqref{asj:eq:sigma1} and \eqref{asj:eq:Dhurwitz}. The Hurwitz-product
class also supplies all principal coefficients and every individual jet.

This is not a theorem that every nonpolynomial arithmetic multiplicity
satisfies those hypotheses. Dirichlet series outside the meromorphic
setting, or with a natural boundary blocking the desired spectral point,
are not covered by the transfer theorem. Their treatment needs additional
structure rather than a formal substitution into its statement.

The previous report already proves the corresponding polynomial and
polynomial-times-logarithm calculus. It also already credits the classical
multiplicative-anomaly mechanism. Those results are antecedents, not new
claims in this package.

\subsection{Corrections and integration safeguards}
No fresh substantive error is asserted in the targeted current source
proofs. The following are proved obstructions to \emph{naive extensions}
and must accompany integration of the new results.

\paragraph{Ordinary derivatives versus Laurent coefficients.}
Equation \eqref{asj:eq:doublepole-example} proves that the double-divisor
shift has residue $-a$ at zero. A formula written as $F_a'(0)$ is
therefore invalid for $a\ne0$; the correct object here is $[s]F_a(s)$.
The coordinate is fixed and must be recorded in any identity store.

\paragraph{The full principal part is necessary.}
The top difference sees only the highest pole, whereas intermediate
local identities involve every $\kappa_{p,h}$. At order $r$ one may
need the numerator through order $r+H$, not merely $r+1$ or $r$.
The exact multiple-pole tests include a case with a third-order pole
at one point and a second-order pole at another.

\paragraph{An Abel limit does not automatically preserve a finite part.}
The divergence in \eqref{asj:eq:twisted-six} and the finite value in
\eqref{asj:eq:untwisted-first-six} prove this explicitly. The compensation
in Theorem~\ref{asj:thm:untwist} is part of the theorem, not a numerical
acceleration that can be dropped.

\paragraph{Do not discard grouped subtractions or base points.}
The Gamma sums and their primitives converge only in their specified
grouped form. The anchored integral vanishes at $a=0$; an arbitrary
negative-polygamma convention cannot replace $K$ without adjusting
its polynomial terms.

\paragraph{Preserve the actual Gaussian status.}
At the inspected revision, the old frozen $S_8$ candidate is rejected,
but a distinct later vector in \texttt{04-S8-candidate.tex} remains
conjectural. Its certified proximity, like the current $S_6$ proximity,
is not a proof. The present arithmetic-spectrum identities neither
prove nor disprove either surviving conjecture. The already proved
$S_4$ identity is not presented as an open target.

\subsection{Reproducible algebra and independent numerical diagnostics}
The accompanying exact program performs \textbf{191} rational or symbolic
polynomial equality checks. They include every coefficient of the
Pochhammer/harmonic identity for $1\le\ell\le18$, the full binomial
recurrence through degree eight, the alternant normalization through
four factors, multiple-pole finite differences, independent Stieltjes
product coefficients, and the six-term normalization with a symbolic
subleading residue. These finite checks do not substitute for the
normal-convergence proofs.

The numerical program uses \textbf{60 decimal working digits} and passes
\textbf{19} equality comparisons. It compares the centered Gamma head
plus an analytically summed Hurwitz tail with an independent zeta-product
series, repeats that comparison for the primitive, and checks reciprocal
polygamma derivatives through order four. It also extracts Laurent
coefficients by a discrete Cauchy integral of the unexpanded meromorphic
binomial series and compares them with the explicit harmonic/Stieltjes
formula through order three. Separate Cauchy extractions check the
six-term double-pole, triple-pole, and polylogarithm-twisted identities.

The largest observed absolute discrepancy in this run was less than
$6.0\times10^{-52}$. The fixed pass threshold was $10^{-35}$; the
smaller observed residual is a diagnostic, not a rigorously enclosed
error. The discrete contour and infinite series were numerically
truncated. The mathematical proofs do not rely on the observed digits.
Complete parameters and individual residuals are in
\texttt{results/numerical\_checks.json}.

A separate limiting diagnostic checks \eqref{asj:eq:first-counter} with
$c=0.9$ and $a=0.12$. It is deliberately not included among the equality
comparisons, because a finite positive $t$ does not give the limiting value.
\begin{center}
\begin{tabular}{@{}cc@{}}\toprule
$t$ & Absolute difference after compensation from the limiting value\\\midrule
$10^{-1}$ & $3.3560095783406\times10^{-2}$\\
$10^{-2}$ & $3.6881187552929\times10^{-3}$\\
$10^{-3}$ & $4.0080828824899\times10^{-4}$\\
$10^{-4}$ & $4.3269361070920\times10^{-5}$\\\bottomrule
\end{tabular}
\end{center}
The proof of the limit is the uniform-contour argument, not the decreasing
sequence of numerical differences.

\subsection{Proposed next research questions}
The following targets are not promoted to theorems in this report.
They are separated from the extension already completed above.
\begin{enumerate}[leftmargin=*,label=\textbf{Q\arabic*.},itemsep=7pt]
\item \textbf{Residue-zero strata for complex twists.}
For products of Lerch factors, classify the twist parameters where the
leading coefficient at $s=1$ vanishes. Determine the resulting exact
degree drops of normalized jets and the transition formulas when several
Lerch factors vanish simultaneously. The first test is one simple zero
of $\Phi(z,1,c)$ with a remaining double Hurwitz pole.

\item \textbf{Several simultaneous untwistings.}
Theorem~\ref{asj:thm:untwist} removes one geometric twist while retaining an
arbitrary Hurwitz product. Construct the finite compensation when two
or more twists approach unity at independent rates. Prove which
rate-dependent finite terms remain after all divergent powers and
logarithms are subtracted. Compare sequential limits with a specified
joint limit rather than silently identifying them.

\item \textbf{Noninteger resonances at other spectral centers.}
At the origin, only integer poles $p=\ell-k$ enter the local formula.
Develop the corresponding finite local calculus about a nonzero total
spectral center so that noninteger or complex poles can resonate.
State the exponent normalization and branch choices before defining
its tangent polynomial law.

\item \textbf{A smaller coordinate algebra for centered Gamma reciprocity.}
For specified rational $a,b,c$, determine which values of $E_{b,c}(a)$
and $I_{b,c}(a)$ reduce to a finite algebra generated by Gamma values,
Hurwitz-zeta derivatives, classical or multiple polylogarithms, and
Dirichlet $L$ derivatives. The comparison algebra must be fixed.
Failure of a symbolic simplifier or a relation search would not prove
irreducibility.

\item \textbf{A reflection or functional equation for the shifted divisor spectrum.}
Construct a useful exact reflection formula in the shift variable for
$F_a(s)$ or its Laurent jets, rather than for the unshifted product
$\zeta(s)^h$. The Gamma representation supplies ordinary derivatives,
but does not by itself give a global reflection law across the negative
spectral lattice.

\item \textbf{Canonical finite-part determinants with several pole orders.}
Our choice $\exp(-[s]F(s))$ is well-defined once the spectral coordinate
is fixed. Determine a normalization with a specified covariance under
spectral rescaling and a specified product law. Record the lower
Laurent coefficients that necessarily enter under a coordinate change;
there is no regulator-independent determinant for free.

\item \textbf{Representation components beyond the Vandermonde.}
For a pole of order $H$, the normalized jet has degree at most $r+H$.
Identify permutation-representation projections whose first occurrence
lies above degree $r$ but below the fully alternating degree. They should
isolate particular Stieltjes residue layers. Give explicit factorizations
and exact finite certificates in the first four-factor cases.

\item \textbf{Other arithmetic products with controlled continuation.}
Apply the transfer theorem to products of primitive Dirichlet $L$
functions and shifted zeta factors, tracking zeros that cancel nominal
poles. Separate residue-class identities from new arithmetic information.
For a base series lacking known continuation, prove that continuation
rather than treating the theorem as evidence for it.

\item \textbf{Formal analytic interfaces and certified regression tests.}
Formalize the finite-principal-part extraction and tangent polarization
independently of the convergence lemmas. Add interval Cauchy extraction
and certified truncation to the current floating-point diagnostics.
The exact finite algebra is already isolated in the supplied program.

\item \textbf{The surviving Gaussian identities remain a separate target.}
An exact certificate for $S_6$, and then the distinct current $S_8$
vector, still needs a relation outside the restricted product-row
presentation identified in the source. The present divisor-spectrum
alternants should not be mistaken for such a certificate.
\end{enumerate}

\subsection{Integration recommendation}
Place the package under a new continuation directory beneath the
Stieltjes-harmonic-resonance report, and mark the explicit Q10 classes
as completed with the hypotheses and scope above. The main candidate
manuscript insertions are the transfer/local calculus, the Gamma
reciprocity with its normalized primitive, and the compensated
polylogarithmic untwisting theorem. The package is self-contained:
its \LaTeX{} build does not depend on unbundled repository files.
No remote repository files were modified while preparing this report.

% END sections/10-audit.tex

\clearpage

% BEGIN references.tex
\begin{thebibliography}{9}
\raggedright
\bibitem{repo}
ProveIt repository, \emph{Polylogarithms research programme},
canonical manuscript and incoming-directory inventory, targeted snapshot
\pin. Relevant chapters: \texttt{07-integration.tex},
\texttt{08-differentiation.tex}, \texttt{10-discovery.tex}, and
\texttt{04-S8-candidate.tex}.
\url{https://github.com/VladimirReshetnikov/ProveIt/tree/3fdd6cc447fac0731801588723109ec3b2e21cec/Analysis/Polylogarithms/docs/manuscript}.

\bibitem{prior}
ProveIt research continuation, \emph{Resonant Gamma Jets and Directional
Zeta Identities: Discrete Stieltjes products, shifted polylogarithms, and
local formulas for higher spectral derivatives}, 10 October 2026.
Especially Section 7, ``Local identities for every derivative of a
spectral product,'' and Section 8, question Q10.
\url{https://github.com/VladimirReshetnikov/ProveIt/tree/3fdd6cc447fac0731801588723109ec3b2e21cec/Analysis/Polylogarithms/docs/reports/stieltjes-harmonic-resonance/continuations/resonant-jets-and-shifts}.

\bibitem{dlmf-hurwitz}
NIST Digital Library of Mathematical Functions,
\emph{Section 25.11: Hurwitz Zeta Function}.
Definitions, continuation, parameter derivatives, and relations to Gamma
and polygamma; accessed during preparation.
\url{https://dlmf.nist.gov/25.11}.

\bibitem{dlmf-gamma}
NIST Digital Library of Mathematical Functions,
\emph{Section 5.7: Series Expansions of the Gamma Function}.
\url{https://dlmf.nist.gov/5.7}.

\bibitem{dlmf-polylog}
NIST Digital Library of Mathematical Functions,
\emph{Section 25.12: Polylogarithms}, especially equation 25.12.12.
\url{https://dlmf.nist.gov/25.12}.

\bibitem{dlmf-lerch}
NIST Digital Library of Mathematical Functions,
\emph{Section 25.14: Lerch's Transcendent}.
\url{https://dlmf.nist.gov/25.14}.

\bibitem{dowker}
J. S. Dowker, \emph{Calculation of the multiplicative anomaly},
arXiv:1412.0549, 2014; corrected version 3, 24 September 2023.
\url{https://arxiv.org/abs/1412.0549}.
\end{thebibliography}

% END references.tex

\end{document}
